%! Boxplots and Comparing Distributions — Unit 1, Lesson 1.6 — Warm-Up
% ONE PAGE at 12pt, blank and key. Three items, no name row (Namestrip).
% GRADUAL-RELEASE lesson: there is no spoiler rule any more, but this page still has no reason to
% name the boxplot vocabulary — it SEEDS the notes rather than previewing them. Item 1 produces
% the five numbers the notes name in their first line (five-number summary); item 2 makes the
% counting in notes problem 13 mechanical; item 3's arithmetic goes on the board and is reused in
% notes problem 20 and again in the AP Practice free response.
% Median, Q1, Q3, IQR, range and the 1.5 x IQR rule were formalized in Lesson 1.5 — prior
% knowledge here.
% PAGE BUDGET: moving 10pt -> 12pt costs roughly a fifth of the vertical budget, so the
% structural spacing was trimmed (itemsep 0.19in -> 0.13in, the leading \vspace, the interior
% \vspace values, and the \blank widths). The \begin{work} block was NOT shaved: it is authored
% BYTE-IDENTICALLY here and in warmup_key/ — under -boxes it ships as a \vphantom (exact height,
% no ink), under -key it prints. That is what keeps the two files one page each.
\documentclass[12pt]{article}
\usepackage{apstats-article}
\usepackage{apstats-boxes}

\begin{document}

\pageheader{Unit 1, Lesson 1.6}{Warm-Up}

\vspace{0.10in}

\begin{enumerate}[leftmargin=1.7em, itemsep=0.13in]

  \item \textbf{From last lesson.} Here are the ages, in years, of nine trees in a city park,
        already put in order.

        \vspace{3pt}
        \begin{center}
        12 \quad 14 \quad 15 \quad 18 \quad 20 \quad 23 \quad 27 \quad 28 \quad 35
        \end{center}
        \vspace{3pt}

        \renewcommand{\arraystretch}{1.3}
        \begin{tabularx}{\linewidth}{@{} Y Y Y Y Y @{}}
        \rowcolor{sky}
        \textbf{Least} & \textbf{$Q_1$} & \textbf{Median} & \textbf{$Q_3$} & \textbf{Greatest} \\
        \hline
        \blank{1.5cm} & \blank{1.5cm} & \blank{1.5cm} & \blank{1.5cm} & \blank{1.5cm} \\
        \end{tabularx}

        \vspace{4pt}
        IQR: \blank{2.2cm} \qquad Range: \blank{2.2cm}

  \item Here are twelve measurements. \textbf{Count} how many fall in each stretch.

        \vspace{3pt}
        \begin{center}
        3 \quad 5 \quad 6 \quad 8 \quad 9 \quad 12 \quad 14 \quad 15 \quad 19 \quad 26 \quad
        33 \quad 55
        \end{center}
        \vspace{3pt}

        \renewcommand{\arraystretch}{1.3}
        \begin{tabularx}{\linewidth}{@{} l Y Y Y @{}}
        \rowcolor{sky}
        \textbf{Stretch} & between 2 and 10 & between 10 and 20 & between 20 and 60 \\
        \hline
        \textbf{How many} & \blank{1.5cm} & \blank{1.5cm} & \blank{1.5cm} \\
        \end{tabularx}

  \item \textbf{From last lesson.} A value counts as unusually far from the rest if it lies more
        than $1.5 \times \text{IQR}$ below $Q_1$ or more than $1.5 \times \text{IQR}$ above $Q_3$.

        \vspace{3pt}
        A data set has $Q_1 = 20$ and $Q_3 = 32$. Is the value \textbf{55} unusually far from the
        rest? Show the arithmetic.

        \begin{work}
        \text{IQR} &= 32 - 20 = 12 \\
        1.5 \times \text{IQR} &= 1.5 \times 12 = 18 \\
        \text{upper boundary} &= Q_3 + 18 = 32 + 18 = 50 \\
        55 &> 50 \quad \Longrightarrow \quad \text{yes, 55 is unusually far}
        \end{work}

\end{enumerate}

\end{document}
